\scanpage{18}
Soit $\omega$ sur $S_{\mathrm{Zar}}$, avec $F_\omega,V_\omega$, on lui associe

$1^\circ)$ $\mathcal M=\varphi^*(\omega),\quad F_{\mathcal M}=\varphi^*(F_\omega),\quad V_{\mathcal M}=\varphi^*(V_\omega)$
\margin{$M=\omega^p$\par
$\varepsilon^*(\Delta(\mathcal M))=\Delta^*(\omega)$\par
\unsure{complexe can. ext.}\par de $\omega^p$ par $\omega^p$}

$2^\circ)$ Triangle
\diagram{\begin{tikzcd}[column sep=small,row sep=large]
&&\mathcal T_2^{(\omega)}=[\omega\overset{V_\omega}{\mathbin{-}}\omega^p]\arrow[dl,"1"']&\\
\unsure{\mathcal T(\omega)|\mathbb T(\omega)}\unsure{:=|:}&\mathcal T_0^{(\omega)}=[\omega^p\xrightarrow{F_\omega}\omega]\arrow[rr]&&\mathcal T_1^{(\omega)}=[\omega^{(p)}\overset0{\mathbin{-}}\omega^{(p)}]\arrow[ul]
\end{tikzcd}}
\unsure{déduit} de $\omega$\quad $(\mathcal T_1^{(\omega)}=\Delta^*(\omega))$

$3^\circ)$ Isom \uuline{can}
\[\varphi^*(\mathcal T(\omega))\overset{\alpha_\omega}\simeq\mathcal T(\varphi^*(\omega))\]

$4^\circ)$ L'\uline{identité}
\[\mathcal T(\omega)_1\overset{\beta_\omega}\simeq\varepsilon^*(\Delta^*(\mathcal M))\simeq[\omega^p\overset0{\mathbin{-}}\omega^p]\]

\uline{Question} Le foncteur
\[
(\omega,F_\omega,V_\omega)\longmapsto(\varphi^*(\omega),\varphi^*(F_\omega),\varphi^*(V_\omega),\mathcal T(\omega),\alpha_\omega,\beta_\omega)
\]
est-il pl. fidèle ?? L'est-il du moins sur la sous-catégorie des triples $(\omega,F_\omega,V_\omega)$ avec $F_\omega=0$ ?
\figure{\begin{tikzpicture}\draw (0,0)--(14,.8);\end{tikzpicture}}

\[
\left\{\begin{aligned}
\varphi^*(\omega)&\xrightarrow{u}\varphi^*(\omega')\\
\mathcal T(\omega)&\overset{\unsure{v|\sigma}}{\mathbin{-}}\mathcal T(\omega')
\end{aligned}\right.
\]
\margin{$u$ compatible avec $F,V$\par
$\unsure{v|\sigma}$ hom de vrais triangles\par
$v_1$ déterminé comme $\varepsilon^*(u)$}
commutativité du diagramme
\diagram{\begin{tikzcd}[column sep=large,row sep=large]
\mathcal T_1(\omega)\arrow[r,"v_1"]\arrow[d,Rightarrow,"\beta_\omega"']&\mathcal T_1(\omega')\arrow[d,Rightarrow,"\beta_{\omega'}"]\\
\varepsilon^*(\Delta^*(\varphi^*(\omega)))\arrow[d,equal]&\varepsilon^*(\Delta^*(\varphi^*(\omega')))\arrow[d,equal]\\
\Delta^*(\omega)\arrow[r,"\varepsilon^*(u)"']&\Delta^*(\omega')
\end{tikzcd}}
\margin{$\varphi^*(\unsure{v|\sigma})$ déterminé par $u$ comme $\mathcal T(u)$}
\diagram{\begin{tikzcd}[column sep=large,row sep=large]
\varphi^*(\mathcal T(\omega))\arrow[r,dash,"\alpha_\omega","\sim"']\arrow[d,"\unsure{\varphi^*(\omega)|\varphi^*(v)}"']&\mathcal T(\varphi^*(\omega))\arrow[d,"\mathcal T(u)"]\\
\varphi^*(\mathcal T(\omega'))\arrow[r,dash,"\alpha_{\omega'}","\sim"']&\mathcal T(\varphi^*(\omega'))
\end{tikzcd}}
\mspage{17}
